% ============================================================
% Documento completo (nuevo) con toda la estructura
% Categorías: Ab, Top, Familias, Mezclas, Functores + IA
% ============================================================

\documentclass[12pt,oneside]{book}

% -------------------- Paquetes básicos -----------------------
\usepackage[spanish]{babel}
\usepackage[T1]{fontenc}
\usepackage[utf8]{inputenc}

\usepackage{amsmath, amssymb, amsthm, mathtools}
\usepackage{geometry}
\geometry{margin=1in}

\usepackage{hyperref}
\hypersetup{
  colorlinks=true,
  linkcolor=blue,
  urlcolor=blue,
  citecolor=blue
}

\usepackage{tikz-cd} % para diagramas con tikz-cd
\usepackage{enumitem}

% -------------------- Entornos -------------------------------
\theoremstyle{definition}
\newtheorem{definition}{Definición}[chapter]
\newtheorem{example}{Ejemplo}[chapter]
\newtheorem{remark}{Observación}[chapter]

% ============================================================
\begin{document}

\frontmatter
\title{Notas: Categorías, Familias, Combinaciones y Aplicación a IA}
\author{Octavio}
\date{\today}
\maketitle

\tableofcontents

\mainmatter

% ============================================================
\chapter{Categorías base}
\label{chap:base}

\section{La categoría \texorpdfstring{$\mathbf{Ab}$}{Ab}: grupos abelianos}
La categoría $\mathbf{Ab}$ se define así:
\begin{itemize}[leftmargin=1.5em]
  \item \textbf{Objetos:} grupos abelianos $(A,+,0,-)$.
  \item \textbf{Morfismos:} homomorfismos de grupos $f:A\to B$.
  \item \textbf{Identidad:} $\mathrm{id}_A:A\to A$.
  \item \textbf{Composición:} si $f:A\to B$ y $g:B\to C$, entonces $g\circ f:A\to C$.
\end{itemize}

\section{La categoría \texorpdfstring{$\mathbf{Top}$}{Top}: espacios topológicos}
La categoría $\mathbf{Top}$ se define así:
\begin{itemize}[leftmargin=1.5em]
  \item \textbf{Objetos:} espacios topológicos $(X,\tau)$.
  \item \textbf{Morfismos:} funciones continuas $g:X\to Y$.
  \item \textbf{Identidad y composición:} las usuales de funciones.
\end{itemize}

% ============================================================
\chapter{Categorías de familias}
\label{chap:familias}

\section{Familias con índice fijo: \texorpdfstring{$\mathcal{C}^I$}{C\^{}I}}
Sea $\mathcal{C}$ una categoría y sea $I$ un conjunto. Considera $I$ como
\emph{categoría discreta} (solo identidades). Entonces:
\[
\mathcal{C}^I \;:=\; \mathrm{Fun}(I_{\mathrm{disc}},\mathcal{C}).
\]
Equivalentemente:
\begin{itemize}[leftmargin=1.5em]
  \item \textbf{Objeto:} una familia $(A_i)_{i\in I}$ con $A_i\in \mathrm{Ob}(\mathcal{C})$.
  \item \textbf{Morfismo:} $(f_i)_{i\in I}:(A_i)\to(B_i)$ con $f_i:A_i\to B_i$ en $\mathcal{C}$.
  \item \textbf{Identidad:} $(\mathrm{id}_{A_i})_{i\in I}$.
  \item \textbf{Composición:} $(g_i)\circ(f_i)=(g_i\circ f_i)_{i\in I}$.
\end{itemize}

\section{Familias con índice variable: \texorpdfstring{$\mathrm{Fam}(\mathcal{C})$}{Fam(C)}}
La categoría $\mathrm{Fam}(\mathcal{C})$ se define así:
\begin{itemize}[leftmargin=1.5em]
  \item \textbf{Objeto:} un par $(I,A)$ donde $I$ es un conjunto y
  $A:I\to \mathrm{Ob}(\mathcal{C})$, es decir, una familia $(A_i)_{i\in I}$ de objetos de $\mathcal{C}$.
  \item \textbf{Morfismo:} de $(I,A)$ a $(J,B)$ es un par
  \[
    (u,(f_i)_{i\in I})
  \]
  donde $u:I\to J$ es una función y, para cada $i\in I$,
  \[
    f_i: A_i \longrightarrow B_{u(i)}
  \]
  es un morfismo en $\mathcal{C}$.
\end{itemize}
La identidad y composición son:
\[
\mathrm{id}_{(I,A)}=\Big(\mathrm{id}_I,\;(\mathrm{id}_{A_i})_{i\in I}\Big),
\]
y si $(u,f):(I,A)\to(J,B)$ y $(v,g):(J,B)\to(K,C)$, entonces
\[
(v,g)\circ(u,f)=\Big(v\circ u,\;(g_{u(i)}\circ f_i)_{i\in I}\Big).
\]

\section{Ejemplo: \texorpdfstring{$\mathrm{Fam}(\mathbf{Ab})$}{Fam(Ab)}}
Un objeto de $\mathrm{Fam}(\mathbf{Ab})$ es una familia de grupos abelianos
$(A_i)_{i\in I}$. Un morfismo incluye una función $u:I\to J$ y homomorfismos
$f_i:A_i\to B_{u(i)}$.

% ============================================================
\chapter{Formas de combinar \texorpdfstring{$\mathbf{Ab}$}{Ab} y \texorpdfstring{$\mathbf{Top}$}{Top}}
\label{chap:combinar}

\section{Producto de categorías: \texorpdfstring{$\mathbf{Ab}\times\mathbf{Top}$}{Ab x Top}}
La categoría producto $\mathbf{Ab}\times\mathbf{Top}$ se define así:
\begin{itemize}[leftmargin=1.5em]
  \item \textbf{Objeto:} un par $(A,X)$ con $A\in\mathbf{Ab}$ y $X\in\mathbf{Top}$.
  \item \textbf{Morfismo:} $(f,g):(A,X)\to(B,Y)$ con $f:A\to B$ en $\mathbf{Ab}$
  y $g:X\to Y$ en $\mathbf{Top}$.
  \item \textbf{Identidad:} $(\mathrm{id}_A,\mathrm{id}_X)$.
  \item \textbf{Composición:} $(f_2,g_2)\circ(f_1,g_1)=(f_2\circ f_1,\;g_2\circ g_1)$.
\end{itemize}
Aquí \emph{no} se exige compatibilidad entre la estructura abeliana y la topológica.

\section{Objetos con ambas estructuras compatibles: \texorpdfstring{$\mathbf{TopAb}$}{TopAb}}
Definimos $\mathbf{TopAb}$ como la categoría de \emph{grupos abelianos topológicos}.
\begin{itemize}[leftmargin=1.5em]
  \item \textbf{Objeto:} un espacio topológico $A$ con mapas continuos
  \[
    +:A\times A\to A,\qquad -:A\to A,\qquad 0:* \to A
  \]
  tales que $(A,+,0,-)$ es un grupo abeliano.
  \item \textbf{Morfismo:} una función continua $f:A\to B$ tal que
  \[
    f(x+y)=f(x)+f(y),\qquad f(0)=0,\qquad f(-x)=-f(x).
  \]
\end{itemize}
Equivalente: $\mathbf{TopAb}$ son los \emph{objetos abelianos internos} en $\mathbf{Top}$.

\section{Plantilla para diseñar una categoría ``mezclada''}
Para construir una categoría de objetos con estructura adicional:
\begin{enumerate}[leftmargin=1.5em]
  \item Elige una categoría ambiente $\mathcal{E}$ (por ejemplo $\mathbf{Top}$).
  \item Especifica los \textbf{datos} (operaciones) como morfismos en $\mathcal{E}$.
  \item Impone los \textbf{axiomas} como diagramas conmutativos (asociatividad, identidad, inversos, conmutatividad, etc.).
  \item Define los \textbf{morfismos} como los que preservan esa estructura.
  \item Verifica identidad y composición: preservan la estructura automáticamente.
\end{enumerate}

% ============================================================
\chapter{Functores entre \texorpdfstring{$\mathbf{Ab}$}{Ab} y \texorpdfstring{$\mathbf{Top}$}{Top}}
\label{chap:funtores}

\section{Qué es un funtor}
Un funtor $F:\mathcal{C}\to\mathcal{D}$ consiste de:
\begin{itemize}[leftmargin=1.5em]
  \item asignación en objetos: $A\mapsto F(A)$,
  \item asignación en morfismos: $f:A\to B \mapsto F(f):F(A)\to F(B)$,
\end{itemize}
tal que:
\[
F(\mathrm{id}_A)=\mathrm{id}_{F(A)},\qquad
F(g\circ f)=F(g)\circ F(f).
\]

\section{Funtor discreto: \texorpdfstring{$D:\mathbf{Ab}\to\mathbf{Top}$}{D:Ab->Top}}
Definimos $D$ por:
\[
D(A)=\text{$A$ con topología discreta},\qquad D(f)=f \text{ (misma función)}.
\]
Entonces $D$ es un funtor porque preserva identidades y composición.

\section{Funtor de olvido: \texorpdfstring{$U:\mathbf{TopAb}\to\mathbf{Ab}$}{U:TopAb->Ab}}
Definimos $U$ por:
\[
U(A,\tau,+,0,-)=(A,+,0,-),\qquad U(f)=f.
\]
Es un funtor que \emph{olvida} la topología.

\section{Homología singular: \texorpdfstring{$H_n:\mathbf{Top}\to\mathbf{Ab}$}{Hn:Top->Ab}}
Para cada $n\ge 0$ existe un funtor (invariante) llamado homología singular:
\[
H_n:\mathbf{Top}\to\mathbf{Ab},
\]
que asigna a un espacio topológico $X$ un grupo abeliano $H_n(X)$,
y a una función continua $g:X\to Y$ un homomorfismo inducido
\[
H_n(g):H_n(X)\to H_n(Y),
\]
cumpliendo:
\[
H_n(\mathrm{id}_X)=\mathrm{id}_{H_n(X)},\qquad
H_n(g\circ f)=H_n(g)\circ H_n(f).
\]

% ============================================================
\chapter{Coproductos y suma directa desde familias}
\label{chap:suma-directa}

\section{Suma directa desde familias en \texorpdfstring{$\mathbf{Ab}$}{Ab}}
Como $\mathbf{Ab}$ tiene coproductos (sumas directas), existe un funtor:
\[
\Sigma:\mathrm{Fam}(\mathbf{Ab})\to\mathbf{Ab},\qquad
\Sigma(I,(A_i))=\bigoplus_{i\in I} A_i.
\]
Este funtor ``colapsa'' una familia de grupos abelianos a su suma directa.

\begin{remark}
En práctica, $\bigoplus_{i\in I} A_i$ es el coproducto en $\mathbf{Ab}$:
elementos con soporte finito sobre el índice $I$.
\end{remark}

% ============================================================
\chapter{Abstracción categórica aplicada a cálculos de IA}
\label{chap:cat-ia}

\section{Motivación: separar estructura de contenido}
En IA es común construir \emph{pipelines} de transformación y modelos compuestos.
La abstracción categórica separa:
\begin{itemize}[leftmargin=1.5em]
  \item \textbf{Estructura:} composición, productos/coprodcutos, invariancias.
  \item \textbf{Contenido:} datos concretos y parámetros entrenables.
\end{itemize}

\section{Ventajas concretas}
\begin{itemize}[leftmargin=1.5em]
  \item \textbf{Composición = pipelines.} Cada etapa es un morfismo; el pipeline es composición.
  \item \textbf{Modularidad.} Cambias un bloque conservando dominio/codominio y todo sigue bien definido.
  \item \textbf{Reuso y pruebas.} Diagramas conmutativos funcionan como pruebas de coherencia.
  \item \textbf{Optimización algebraica.} Reasociar/fusionar cómputos sin cambiar resultados.
  \item \textbf{Batching como familias.} Un mini-batch se modela en $\mathcal{C}^I$.
  \item \textbf{Invariancias explícitas.} Simetrías como acciones de grupo reducen complejidad efectiva.
  \item \textbf{Interoperabilidad.} Functores conectan representaciones (imagen$\to$grafo$\to$TDA$\to$vector).
\end{itemize}

\section{Ejemplo 1: Pipeline de IA como composición}
Define una categoría (conceptual) $\mathcal{P}$:
\begin{itemize}[leftmargin=1.5em]
  \item \textbf{Objetos:} espacios de datos (Raw, Tokens, Embeddings, Logits, \dots).
  \item \textbf{Morfismos:} transformaciones computables (tokenizar, normalizar, embed, clasificar).
\end{itemize}

Un pipeline típico:
\[
\mathrm{Raw}
\xrightarrow{\;\mathrm{tokenize}\;}
\mathrm{Tokens}
\xrightarrow{\;\mathrm{embed}\;}
\mathrm{Emb}
\xrightarrow{\;\mathrm{head}\;}
\mathrm{Logits}.
\]
El modelo completo es el morfismo compuesto:
\[
\mathrm{head}\circ \mathrm{embed}\circ \mathrm{tokenize}:\mathrm{Raw}\to \mathrm{Logits}.
\]

\section{Ejemplo 2: Mini-batching como familia}
Sea $\mathcal{C}$ una categoría (p.ej.\ $\mathbf{Vect}$) e $I$ un conjunto índice.
Entonces:
\[
\mathcal{C}^I=\mathrm{Fun}(I_{\mathrm{disc}},\mathcal{C}).
\]
Un batch de embeddings:
\[
(E_i)_{i\in I}\in \mathbf{Vect}^I.
\]
Aplicar una transformación por componente:
\[
(f_i:E_i\to E_i')_{i\in I}:(E_i)_{i\in I}\to (E_i')_{i\in I},
\]
y composición por componente:
\[
(g_i)\circ(f_i)=(g_i\circ f_i)_{i\in I}.
\]

\section{Ejemplo 3: TDA como cadena de traductores (funtores)}
Esquema conceptual:
\[
\text{Datos}
\longrightarrow
\text{Complejos simpliciales}
\longrightarrow
\text{Módulos de persistencia}
\longrightarrow
\mathbf{Vect}.
\]
Más explícito:
\[
(X,d)
\xmapsto{\;\mathrm{VR}_\bullet\;}
\{\mathrm{VR}_\epsilon(X)\}_{\epsilon\ge 0}
\xmapsto{\;H_k\;}
\{H_k(\mathrm{VR}_\epsilon(X))\}_{\epsilon\ge 0}
\xmapsto{\;\Phi\;}
\mathbb{R}^m.
\]

\section{Ejemplo 4: Backprop y composición (lectura estructural)}
En entrenamiento:
\[
\theta \leftarrow \theta - \eta \nabla_{\theta}\mathcal{L}.
\]
Si el modelo es una composición:
\[
f=f_n\circ f_{n-1}\circ\cdots\circ f_1,
\]
la regla de la cadena expresa compatibilidad con la composición:
\[
D(f_n\circ\cdots\circ f_1)=Df_n\circ\cdots\circ Df_1,
\]
con las composiciones adecuadas (jacobianos/diferenciales).

\section{Reglas prácticas: tamaño y control de combinaciones}
No hay un límite universal de combinaciones, pero sí restricciones por \emph{tamaño}.

\subsection{Categorías grandes y familias pequeñas}
Permitir índices arbitrarios en $\mathrm{Fam}(\mathcal{C})$ puede producir una categoría
muy grande. Para controlar esto se usan restricciones típicas:
\begin{itemize}[leftmargin=1.5em]
  \item $\mathrm{Fam}_{\mathrm{fin}}(\mathcal{C})$: solo familias con índice finito.
  \item $\mathrm{Fam}_{\mathbb{N}}(\mathcal{C})$: solo familias numerables.
  \item $\mathrm{Fam}_{U}(\mathcal{C})$: índices ``pequeños'' respecto a un universo de Grothendieck $U$.
\end{itemize}

% ============================================================
\backmatter
\chapter*{Notas finales}
Si quieres, en el siguiente paso podemos agregar:
\begin{itemize}[leftmargin=1.5em]
  \item diagramas conmutativos de axiomas de grupo en $\mathbf{TopAb}$ usando \texttt{tikz-cd},
  \item un ejemplo completo de acción de grupo e invariancia (p.ej.\ traslaciones en imágenes),
  \item una sección de ``categorías para deep learning'' (Monoidal categories / string diagrams).
\end{itemize}

\end{document}
